\documentclass[11pt,a4paper]{article}
\usepackage[margin=20mm,top=27mm,bottom=20mm,headheight=25pt,footskip=8mm]{geometry}
\usepackage[utf8]{inputenc}
\usepackage[T1]{fontenc}
\usepackage[brazil]{babel}
\usepackage{helvet,amsmath,amssymb,booktabs,array,tabularx,enumitem,graphicx,xcolor,tcolorbox,fancyhdr,hyperref,tikz}
\usetikzlibrary{arrows.meta,positioning}
\renewcommand{\familydefault}{\sfdefault}
\definecolor{azul}{RGB}{0,83,134}
\definecolor{ouro}{RGB}{206,177,68}
\hypersetup{pdflang=pt-BR,colorlinks=true,urlcolor=azul,linkcolor=azul,pdfauthor={Luan F. Nogueira e Luiz F. de Melo Soares}}
\setlength{\parindent}{0pt}
\setlength{\parskip}{5pt}
\setlist{nosep,leftmargin=*}
\pagestyle{fancy}\fancyhf{}
\fancyhead[L]{\small\color{azul}DO TRANSISTOR AO CHIP}
\fancyhead[R]{\includegraphics[height=7mm]{../slides/libs/emblemufc.pdf}}
\fancyfoot[L]{\scriptsize Luan F. Nogueira • Luiz F. de Melo Soares}
\fancyfoot[R]{\scriptsize\thepage}
\renewcommand{\headrulewidth}{.4pt}
\setlength{\headheight}{25pt}
\newcommand{\titulo}[2]{{\LARGE\color{azul}\bfseries #1}\par{\large #2}\par\vspace{2mm}}
\newcommand{\caixa}[2]{\begin{tcolorbox}[colback=azul!4,colframe=azul,title={#1},fonttitle=\bfseries,boxrule=.5pt,arc=1mm]#2\end{tcolorbox}}
\newcommand{\VDD}{V_{\mathrm{DD}}}

\hypersetup{pdftitle={Guia dos facilitadores — oficinas visuais de duas horas}}
\begin{document}
\titulo{Guia dos facilitadores}{Três oficinas de duas horas · versão visual e participativa}
\textbf{Proposta.} Usar a figura como ponto de partida e pedir uma previsão antes de avançar. Os quadros mostram os conceitos essenciais; estas notas e os gabaritos guardam o aprofundamento. As aulas 1 e 2 têm atividades dentro dos 120 minutos. A prática no Virtuoso da aula 3 será elaborada posteriormente.

\textbf{Como apresentar.} Abra o PDF principal em modo de apresentação e avance pelas páginas. Um quadro pode ter várias etapas: o número no rodapé permanece o mesmo durante a sequência. O PDF \texttt{resumo} tem uma página por quadro, selecionando a etapa de síntese. Os PDFs \texttt{notas} acrescentam orientações para quem apresenta.

\textbf{Aula 1 · do silício ao MOSFET.}
\begin{center}\small\begin{tabularx}{\linewidth}{lX}\toprule Minutos & Condução\\\midrule
0--10 & Apresentação dos facilitadores e percurso das três oficinas.\\
10--30 & Materiais, ligações, elétron/lacuna e dopagem.\\
30--50 & Junção PN: prever difusão, observar depleção e comparar polarizações.\\
50--60 & Pausa.\\
60--90 & Formar o canal, comparar NMOS/PMOS e resolver o exemplo de tensões.\\
90--105 & Inversor: rastrear carga e descarga; distinguir esquema de layout.\\
105--120 & Desenho em dupla, explicação e fechamento.\\\bottomrule\end{tabularx}\end{center}
\textbf{Aula 2 · do transistor às portas.}
\begin{center}\small\begin{tabularx}{\linewidth}{lX}\toprule Minutos & Condução\\\midrule
0--10 & Retomar as chaves e a interpretação dos níveis de tensão.\\
10--30 & Inversor: forma de onda, carga, atraso e energia.\\
30--50 & Redes série/paralelo, construção da NAND e quatro estados.\\
50--60 & Pausa.\\
60--80 & Construção da NOR, quatro estados e comparação das portas.\\
80--105 & Desenho, troca de circuitos entre duplas e extensão a três entradas.\\
105--120 & Composição de portas, previsão para a prática e fechamento.\\\bottomrule\end{tabularx}\end{center}
\textbf{Ritmo.} Tratar cada clique como uma pergunta curta: ``o que mudou?'', ``qual caminho fechou?'' ou ``de onde vem a carga?''. Usar o tempo disponível para ouvir as respostas e desenhar junto. Uma pessoa pode apresentar enquanto a outra acompanha as duplas e recolhe dúvidas.
\newpage
\titulo{O que observar nas figuras}{Roteiro das sequências e limites dos modelos}
\begin{description}[style=nextline]
\item[Deslocamento de lacuna] Identificar onde fica o estado vazio depois que um elétron o preenche. O desenho mostra ocupação de estados, não órbitas atômicas.
\item[Depleção da junção PN] Começar com os portadores majoritários; prever a difusão; revelar a região com íons fixos. A região é pobre em portadores móveis, mas contém carga. O campo interno aponta de N para P.
\item[Polarização da junção] Comparar largura da depleção e barreira. Direta: P positivo em relação a N. Reversa: P negativo. As setas de campo mostram orientação; não quantificam seu módulo.
\item[Ativação do NMOS] Identificar dielétrico, depleção e inversão. A porta controla por campo; os portadores do canal não atravessam o dielétrico. Com dreno positivo, corrente convencional e elétrons se deslocam em sentidos opostos.
\item[Inversor] Rastrear o caminho verde até um trilho e apontar a chave cinza que interrompe o outro. Os estados mostrados são estáveis e idealizados.
\item[NAND/NOR] A sequência 00, 01, 11, 10 muda um bit por vez. Verde marca uma chave habilitada no modelo ideal; uma rede só conduz se existir caminho completo. Revelar as tabelas depois das previsões.
\end{description}
\textbf{Critérios de aprendizagem.} O participante deve distinguir porta de saída, calcular diferenças de tensão entre terminais e justificar tabelas pela conectividade. A saída flutuante não equivale a zero. Baixo consumo estático ideal não equivale a energia de comutação nula.

\textbf{Limites didáticos.} Modelos de enriquecimento, corte planar conceitual e CMOS estático complementar. Os corpos NMOS ficam em GND e os corpos PMOS em $\VDD$ no modelo bulk. Os limiares constantes de 0,4 V e a alimentação de 1,2 V são exemplos; não definem PDK ou parâmetros da prática. Verde é uma abstração de habilitação: nós internos de pilhas exigem análise elétrica para prever correntes reais.

\textbf{Sobre carga e energia.} Os equivalentes RC e as formas de onda são qualitativos. A energia se dissipa nos caminhos de carga e descarga; a capacitância armazena energia temporariamente. Fugas e corrente direta transitória são efeitos adicionais. Caso haja interesse, a estimativa capacitiva é $P\approx\alpha C_L\VDD^2f$, com $\alpha$ contando eventos 0$\rightarrow$1 por ciclo.
\newpage
\titulo{Encerramento da aula 3}{Trilhas de estudo e âmbito Brasil}
\textbf{Duração.} Reservar cerca de 25 minutos finais para este bloco e 95 minutos para a futura prática. O roteiro do Virtuoso continua pendente; esta entrega não define comandos, PDK, dimensões ou condições de simulação.

\begin{center}\small\begin{tabularx}{\linewidth}{lX}\toprule Minutos do fechamento & Uso sugerido\\\midrule
0--10 & Cinco trilhas e exemplo de produto: uma pergunta e um projeto por área.\\
10--18 & Brasil: situar os polos e destacar exemplos da cadeia.\\
18--25 & Escolha de um primeiro projeto, conversa e fontes para consulta.\\\bottomrule\end{tabularx}\end{center}
\textbf{Como usar o mapa.} A visão nacional abre a seção; as seis páginas regionais preservam a relação integral da planilha. Elas são material de consulta: escolher um ou dois exemplos por grupo e relacionar a atuação ao que a turma aprendeu. O mapa A3 pode ser distribuído como uma folha de referência.

\textbf{Créditos Brasil.} Levantamento de Matheus Grossi, \href{https://github.com/MattGrossi12/brazil_silicon_roadmap}{Brazil Silicon Roadmap}, usando a planilha PT-BR fornecida em 28/10/2026. Esta edição contém 54 organizações, 73 presenças por polo e 13 polos; um nome pode aparecer em vários locais. Empresas são distinguidas de centros e grupos de pesquisa. A relação abrange projeto, EDA, fabricação, integração, eletrônica aplicada e pesquisa; a atuação detalhada e as incertezas da fonte estão nos CSVs.

\textbf{Geografia.} Contorno simplificado de \href{https://github.com/johan/world.geo.json/blob/master/countries/BRA.geo.json}{johan/world.geo.json}; desenho recomposto em TikZ. Pontos indicam polos, não endereços. Regiões descritas como interior são áreas agregadas. A imagem do roadmap fornecida pelo usuário não é colada na apresentação.

\textbf{Mini-fichas.} Luan F. Nogueira tem percurso na UFC e na Polytech Montpellier e experiência em sistemas embarcados, FPGA e RTL de aceleradores no imec. Luiz F. de Melo Soares tem percurso na UFC e no ISMIN/Mines Saint-Étienne e experiência em projeto analógico CMOS, simulação e layout no CPPM. As descrições se baseiam nos currículos fornecidos, sem presumir conclusão de curso ou vínculo atual.

\textbf{Autoria e fontes.} Minicurso: Luan F. Nogueira e Luiz F. de Melo Soares. Template Beamer UFC: Maurício Moreira Neto, classe original 4:3. Referências técnicas: Fathi Salem, J.-M. Dutertre e Steve Ward. Figuras e sequências: elaboração didática do minicurso. O Git local guarda fontes, figuras necessárias e dados derivados; referências originais, currículos e arquivos gerados ficam fora do versionamento.
\end{document}
